\section{From a Case III residual chain to automorphy}

A residual chain joins representations in characteristic $\ell$. We need
an ideal that vanishes at the original integral representation. The two
ends of the chain may lie on intersections of components. We first move
support across these intersections, then lift it to the original point.

% BEGIN CJP1
\paragraph{CJP1. Moving support backwards along a rank-four chain.}
Let $\Lambda$ be the complete coefficient DVR, $k$ its finite residue
field, $X$ the space for the standard frame, $B=X.B$, $A=X.A$ and
$J\subset B$. At a residual point $x:A\to k$, \emph{some} means that
some minimal prime below $\ker x$ contains $JA$; \emph{all} means
every such minimal prime does. Assume the original implication from
some to all at residual $\mathrm{Good4}$ points. Suppose the pair at
$x_0$ is residual, a $\mathrm{Chain4}$ joins the pairs at $x_0,x_1$
in that order, and some holds at $x_1$. Then all holds at $x_0$.

Indeed, each polynomial link is represented by $b:A\to k[U]$ by the
line property of $X$. Evaluating at $0,1$ gives the endpoint pairs.
It also gives the specified endpoint maps whenever those have already
been chosen. Both maps agree on $B$ with its residue map: any
$\Lambda$-algebra map $B\to k$ is surjective, has the unique maximal
ideal as kernel, and every residue has a representative in $\Lambda$.
The other algebra generators are entries of the representation and of
its local matrices, on which equality of the pairs gives equality.
Both endpoint kernels thus contain $\ker b$. Choose a minimal prime
below this prime. All at the right endpoint implies that this minimal
prime contains $JA$, witnessing some at the left. Evaluation gives a
residual point even when only its pair was specified. The initial pair
is good by the chain definition, and every later endpoint is good by
the link clause and residual by evaluation. Upgrade some to all at the
terminal point, then move backwards one link at a time, upgrading at each
vertex. At length zero the generator argument identifies the endpoint
maps. Goodness at the interior points of a polynomial line is unnecessary.
% END CJP1

% BEGIN CJP2
\paragraph{CJP2. Lifting support from the seed to the prescribed valued point.}
Fix a fit, deep tower $T$, an embedding $\iota$, and the small pair
$P$ over the ambient valuation ring $V$. It is a pair for the empty
auxiliary frame and its reduction is residual. Let $Q$ over $O_1$
be a seed pair at $\iota$. Suppose a rank-four chain over $k_V$ joins
the reduction of $P$ to that of $Q$. Assume cohomology,
$\mathrm{PMExAlt}$ and $\mathrm{ChainRankFour}$, with frame bound
$C$ for every residual coefficient map and $d(T)\ge30(C+1)^2$.

RCD3 supplies a coefficient threshold for the entire chain. The small-pair
construction of UX11 supplies a second, and $\mathrm{PMExAlt}(T,\iota)$
a third, before $X$ is chosen. Take a common upper coefficient DVR
$\Lambda$ for these three rings and $O_1$ by successive composita. Its
residue $k$ is finite over that of $\mathbb Z_\ell$. Standardness of
the empty frame on a fit tower supplies the space $X$ of UX6. Apply
the tower supplier to this $X$ for its ideal $J\subset B$.
Write $j:k\hookrightarrow k_V$ for the residue inclusion. RCD3 gives
endpoints $s,t$ and their chain over $k$ with the prescribed images
under $j$. UX11 gives a complete local $\Lambda$-algebra $R$ with
residue $k$, an injection $f:R\to V$ and a continuous pair $P_R$
mapping to $P$, with
\[
 f|_\Lambda=a_\Lambda,\qquad
 \operatorname{res}_V f=j\operatorname{res}_R .
\]
The pair equations descend along $f$ by injectivity. The reductions of
$P_R$ and $Q_\Lambda$ equal $s,t$ since their images by $j$ agree. The
coefficient inclusion is local: a unit upstairs maps to a unit in $V$,
whose coefficient maps are local. Continuity therefore persists for the
seed. The two applications of universality give $\Lambda$-algebra maps
\[
 f_0:A\longrightarrow R,\qquad f_1:A\longrightarrow\Lambda ,
 \qquad X.P_{f_0}=P_R,\quad X.P_{f_1}=Q_\Lambda .
\]
Their residual maps $x_0,x_1$ realize $s,t$. The composite of $f_1$
with $\Lambda\hookrightarrow\overline{\mathbb Q}_\ell$ gives the
original seed, hence $\mathrm{LocZero}$. Choose a minimal prime below
the kernel of this composite. Its contraction to $B$ is killed by the
seed and hence contains $J$. The composite kernel lies in $\ker x_1$
by injectivity of $\Lambda\hookrightarrow\overline{\mathbb Q}_\ell$.
This gives some at $x_1$. The map to $k$ kills the original maximal
ideal by locality. The bound and degree now give the implication
$\mathrm{ChainRankFour}$ for this $X,J$. CJP1 gives all at $x_0$.

Put $y=ff_0$. Its kernel is prime since $V$ is a domain and lies in
$\ker x_0$ by injectivity of $f$. A minimal prime below $\ker y$
therefore contains $JA$. Thus $y(JA)=0$. On $B$ the reduction of
$f_0$ is the canonical residue map. The residue compatibility above
gives $y(\mathfrak m_B)\subseteq\mathfrak m_V$. We have
$y|_\Lambda=a_\Lambda$ and $X.P_y=P$ by construction.
% END CJP2

% BEGIN CJP3
\paragraph{CJP3. The Case III chain-to-promodularity implication.}
Assume the original tensor-factor tower, rank-four component implication
and seeded rank-four chain reduction. Use RCD3 for chain descent. Under
cohomology, Ferrero--Washington, $\mathrm{HFit}$, $\mathrm{SeedIVAlt}$
and $\mathrm{PMExAlt}$, every Case III star pair admits
$\mathrm{ProModAlt}$ at its prescribed $\iota$.

Fix $D,\iota$, the ambient ring $V$ with fraction field $K$, and the
given pair $(r,u)$. The first column of the residual ordinary flag picks
a permitted label: in the nonidentity permutation case its eigenvector
equation gives the required local character equality (UX2). Apply the
tensor-factor tower hypothesis to this label to get $T_0$ satisfying
the original acceptability, four-factor and tameness conditions.
$\mathrm{HFit}$ gives fitness, so the distinguished element separates
the four residual diagonal characters. At this element the label
construction gives a near-split integral pair conjugate to $r$ over
$K$. The lattice construction of UX2 applied to the star alternative
gives a further integral pair $P'$ in the prescribed lattice shape.
Its reduction is the input of the seeded rank-four reduction at
$T_0,\iota,V$. The resulting tame extension $T_1$ works for this
reduction and all later deep towers.

Fitness passes to the extension. Choose $C$ at $T_1$ from the uniform
cohomological and Ferrero--Washington bounds of UX5, then take a deep
cyclotomic layer $T$ with $d(T)\ge30(C+1)^2$. The reduction gives a
seed pair $Q$ on $T$ and a chain for every residual pair lying over
the specified global reduction. Restrict $P'$ using the places and
flags of UX11. This gives a small pair $P$ with $P.r=P'.r|_{G_F}$,
with residual pair over the specified global reduction, and with
\[
 P.\psi_{v,i}(s)=a_O\bigl(D.\mathrm{ps}_i(\mathrm{res}_{\mathbb Q}s)\bigr)
       \qquad (s\in I_{F,v}).
\]
This is an eligible starting point for the chain to $Q$. CJP2 supplies
$X,J,y$ with $X.P_y=P$. Take the global representation $r_1=P'.r$ over
$V$. The two conjugacies compose over $K$ by multiplication of the
conjugating matrices; the multiplier is $a_O(D.\vartheta)$ by the
global pair condition. The restriction of $r_1$ is $(X.P_y).r$.
Functoriality of the local character under $y$ gives the required
inertia equalities for $X.P$ from the display. Along with fitness,
depth, the tower record, coefficient compatibility, locality and
annihilation of $J$, these are all fields of $\mathrm{ProModAlt}$.
% END CJP3

% BEGIN CJA1
\paragraph{CJA1. One embedding and the two host cases.}
Fix the inputs of the source rank-four density statement: the original
host $h$ with $h.\mathrm{Big}(t)$, an odd prime $\ell$, its coefficient
datum, $0<e<t$, $\psi\in h.X(c,e)$ and a good $D$ with
$\mathrm{OfCase}(D,h,e,\psi)$. Fix an isomorphism
$\iota:\overline{\mathbb Q}_\ell\simeq\mathbb C$, available from the
characteristic-zero ambient datum. Assume the source cohomology,
Ferrero--Washington, fitness, both seed inputs, and the shared
potential-modularity tower supplier. Suppose Case III promodularity is
given in its original form. Then for every ambient field and every
$\mathrm{KPair}$ $(r,u)$ with a star representation over its fraction
field, there is an actual $\mathrm{ProModAlt}(D,\iota,\mathfrak a,r)$.

To see this, expand $\mathrm{OfCase}$. In Case II its character data
have parameters $(n_{\rm II},1,1)$, and in Case III they have parameters
$(n_{\rm III},p_{\rm One}(h.\chi),\operatorname{ord}(\psi))$. Together
with goodness, the displayed host, exponent and character supply the
correct context in each branch. In Case II apply UX12 with the
rank-two seed input and shared tower supplier. In Case III apply the
given Case III implication with cohomology, Ferrero--Washington,
fitness, the rank-four seed and the same shared supplier, then evaluate
it on these very $D,\iota,\mathfrak a,r,u$. The resulting tower, ideal
and specialization may depend on the pair. The embedding was chosen
before any pair or ambient. CJP3 supplies this Case III implication
when its stated rank-four inputs are available.
% END CJA1

% BEGIN CJA2
\paragraph{CJA2. Finiteness, the block obstruction and automorphy on the original datum.}
At the common $\iota$ of CJA1 the source concludes
\[
 D.\mathrm{Fin0}\ \land\ D.\mathrm{NoEnd}\ \land\
 D.\mathrm{AllIrr}(\mathrm{IsAut4Alt}\,\iota).
\]
Here finiteness concerns the original weight-zero monomial algebra of
every allowed lift. The second clause excludes full span of the designated
block at a block point. The third concerns every irreducible algebraic
point of the same $D$.

For a characteristic-$\ell$ star pair use its witness from CJA1. The
finite special weight fibre makes the characteristic polynomial on
the tower subgroup equal to the product of the four residual character
factors. Taking powers places every eigenvalue over a finite field;
the difference of the two traces lies in the maximal ideal, so is zero.
The resulting commutator obstruction of IX9 excludes both full
span and the nonsplit block extension. This gives $\mathrm{Con}(D)$.
In characteristic zero take the prescribed coefficient identification
$K\simeq\overline{\mathbb Q}_\ell$. For every star pair, IX10 transports
the purity clause along it and along the witness conjugacy. For the
resulting representation there are an open subgroup and a finite
exceptional set such that at a permitted Frobenius of degree $f>0$
its roots satisfy $|\iota x|<p^f|\iota y|$ whenever $y\ne0$.
This proves $\mathrm{PurityAtPoints}$, including the extension
alternative. On irreducible pairs use full span to supply the first
star alternative. The automorphy clause and its transport in IX10 give
the source automorphic attachment $\mathrm{IsAut4Alt}$ for the same
$\iota$. This gives $\mathrm{AutmAlt}$ on its full domain.

Apply IX8 to these three assertions. In its finiteness argument a
characteristic-$\ell$ curve would either violate the star obstruction or
make the selected weight-zero monomial periodic. In its block argument
the integral nonsplit construction and the root inequality exclude the
proposed end point. Its irreducible-point argument uses the just-proved
automorphic attachment. We obtain all three conclusions, and may now
introduce the single existential embedding in the source density
statement. This proves the automorphy implication on its stated inputs.
The tensor-factor tower, rank-four component and seed hypotheses retain
their separate mathematical content.
% END CJA2
